\newpage
\section{Extended Literature Review}
\label{sec:lit-review}

Looping model depth has been well explored by prior work; with a large body of work studying looping within general language modeling \citep{dehghaniUniversalTransformers2019, zhuScalingLatentReasoning2025, geiping_scaling_2025, mcleish_retrofitted_recurrence, baeMixtureofRecursionsLearningDynamic2025} or small-scale algorithmic problems \citep{avi_learn_algorithm, yangLoopedTransformersAre2023, bansalEndtoendAlgorithmSynthesis2022, wangHierarchicalReasoningModel2025b, jolicoeur-martineauLessMoreRecursive2025}. Within looped architectures, the design of training paradigms can be relatively split between architectures with explicit halting mechanisms \citep{dehghaniUniversalTransformers2019, zhuScalingLatentReasoning2025, baeMixtureofRecursionsLearningDynamic2025, jolicoeur-martineauLessMoreRecursive2025, wangHierarchicalReasoningModel2025b} and those with implicit halting mechanisms \citep{geiping_scaling_2025, mcleish_retrofitted_recurrence, LoopFormerElasticDepthLooped2025, xuExpressivePowerLooped2025}. 
Looped architectures trained with an explicit halting mechanism use specialized architectures to predict when to early exit tokens, preventing additional computation updates on their recurrent stream \citep{wangHierarchicalReasoningModel2025b, jolicoeur-martineauLessMoreRecursive2025, baeMixtureofRecursionsLearningDynamic2025, dehghaniUniversalTransformers2019, elbayadDepthAdaptiveTransformer2020}. Specifically, \citet{wangHierarchicalReasoningModel2025b, jolicoeur-martineauLessMoreRecursive2025} formalize \emph{adaptive-computation-time}, a method that utilizes Q-learning as a means to determine convergence. Similarly, works such as \citet{baeMixtureofRecursionsLearningDynamic2025} define an architecture that uses light-weight routers to assign dynamic recursion depths, while \citet{zhuScalingLatentReasoning2025} uses a prediction head to dynamically define a probability of exiting after recurrent passes. A majority of these approaches draw on methods of layer skipping \citep{elhoushiLayerSkipEnablingEarly2024, raposoMixtureofDepthsDynamicallyAllocating2024}; however, these methods differ from using a shared parameterization for a recurrent block.

Alternatively, looped architectures with an implicit halting mechanism, such as \citet{geiping_scaling_2025, mcleish_retrofitted_recurrence,avi_learn_algorithm,bansalEndtoendAlgorithmSynthesis2022}, train models with stochastically sampled recurrent steps during pretraining, and then use the KL-divergence between two successive steps to decide when to exit from the recurrent unit early. Finally, \citet{LoopFormerElasticDepthLooped2025} ignores adaptive early exiting altogether, instead pretraining a recurrent unit on a static number of recurrences and enforcing a consistency loss on intermediate recurrences.
Our work focuses solely on implicit recurrent depth models \citep{geiping_scaling_2025, mcleish_retrofitted_recurrence}, which are derived from prior initial work \citep{avi_learn_algorithm, bansalEndtoendAlgorithmSynthesis2022}.

Beyond training paradigms, there are several differing architectural design choices for looped models \citep{geiping_scaling_2025,bansalEndtoendAlgorithmSynthesis2022, saunshiReasoningLatentThoughts2025b}. In simple looped architectures that only place a single recurrent unit, the placement of the looped unit is non-trivial, with certain works looping over all layers \citep{dehghaniUniversalTransformers2019,Csordas2024MoEUTMU,Bae2024RelaxedRT}. Alternatively, \citet{saunshiReasoningLatentThoughts2025b} find middle-looping recurrent units are the most effective in comparison to other formulations, such as pre-looping and post-looping, which loop the beginning and end of the model. The effectiveness of Middle-looping is consistent with the initial work in synthetic problems by \citet{bansalEndtoendAlgorithmSynthesis2022, avi_learn_algorithm} and with the architecture choices of \citet{geiping_scaling_2025,mcleish_retrofitted_recurrence} in large-scale language models before training. Within middle-looping architectures, the number of layers within each unit is mostly chosen ad hoc; however, when bootstrapping from a baseline model, \citet{koishekenov2025encodethinkdecodescaling} found that you optimize placement by algorithmically selecting layers within a model to loop.

While these prior formulations of looping focus on a single recurrent block, hierarchical \citep{wangHierarchicalReasoningModel2025b,jolicoeur-martineauLessMoreRecursive2025}, parallel \citep{wu2025parallellooptransformerefficient}, and multi-step \citep{jacobs2026blockrecurrentdynamicsvisiontransformers} formulations of layer looping exist. Furthermore, while not all under the same architectural paradigm, layer looping has been explored in multiple domains (e.g., language \citep{geiping_scaling_2025,mcleish_retrofitted_recurrence}, images \citep{jacobs2026blockrecurrentdynamicsvisiontransformers}, multi-modal systems \citep{alabdulmohsin2025recursiveinferencescalingwinning}, synthetic algorithmic problems \citep{avi_learn_algorithm,bansalEndtoendAlgorithmSynthesis2022,yangLoopedTransformersAre2023}), with the choice of looping style and model architecture design changing based on the specific modality. Where layer looping is introduced, how it is affected by individual modalities, and efficient, FLOP-optimal implementations of layer looping remain open questions.

Finally, layer looping is often deeply tied to deep equilibrium (DEQ) models \citep{bai2019deepequilibriummodels,bai2022neural}, due to the fixed-point nature often learned in recurrence. DEQs find the equilibrium points via root-finding to approximate an \emph{infinite depth} network. However, unlike looped architectures trained with truncated backpropagation, a key advantage of DEQ models is their use of implicit differentiation through \emph{infinite depth}, which keeps memory constant and independent of effective depth used to solve the fixed point using a rooting finding algorithm. While the use of implicit differentiation in DEQs enables more efficient training, we focus on work that does explicit backpropagation rollouts \citep{geiping_scaling_2025, mcleish_retrofitted_recurrence, bansalEndtoendAlgorithmSynthesis2022, yang2024loopedtransformersbetterlearning}.
Within looped architectures, \citet{geiping_scaling_2025,mcleish_retrofitted_recurrence} adopt the usage of path independence from equilibrium models \citep{anilPathIndependentEquilibrium} to warrant their choice of $h_0$ initialization.
